\section{Variantes exploratoires BGL, conçues après lecture du test}
\label{sec:annexe-variantes-bgl}

Les variantes ci-dessous complètent le protocole principal BGL (section~\ref{sec:bgl-resultats}). Elles ont été explorées sans plan préalable, après observation du test : elles décrivent des mécanismes et ne constituent pas des estimations indépendantes de la performance. Toutes utilisent la même plage de test, sauf la partition chronologique 60/20/20, un Drain3 ajusté sans label sur le seul entraînement puis gelé, des histogrammes calculés sur les événements normaux de l'entraînement et un seuil choisi sur la validation.

\begin{table}[H]
\centering
\caption[BGL : variantes exploratoires.]{BGL : variantes exploratoires. Le seuil est choisi sur la validation, mais le choix des variantes dépend de l'observation du test. Source : campagnes BGL, second protocole et protocole contigu (tableau~\ref{tab:preuves-experimentales}).}
\label{tab:bgl-protocoles}
\small
\resizebox{\textwidth}{!}{%
\begin{tabular}{lrrrr}
\toprule
\textbf{Protocole} & \textbf{N test / positifs} & \textbf{Positifs connus} & \textbf{F1 Histogram} & \textbf{F1 template inconnu} \\
\midrule
Principal ($49\,999$ lignes) & $19\,908$ / $2\,885$ & 0 & 0,9137 & 0,2779 \\
5 blocs de $10\,000$ lignes & $19\,908$ / $2\,885$ & 4 & 0,9238 & 0,2801 \\
10 blocs & $19\,908$ / $2\,885$ & 5 & 0,8913 & 0,2803 \\
20 blocs & $19\,908$ / $2\,885$ & 6 & 0,7309 & 0,2924 \\
Drain3 adapté sur validation & $19\,908$ / $2\,885$ & 4 & 0,9137 & 0,2796 \\
Chronologique 60/20/20 & $97\,107$ / $2\,439$ & 7 & 0,5985 & 0,0576 \\
20 blocs avant le test & $19\,908$ / $2\,885$ & 6 & 0,9405 & 0,9156 \\
G : $80\,000$ lignes contiguës & $19\,908$ / $2\,885$ & $2\,880$ & 0,9419 & 0,0033 \\
\bottomrule
\end{tabular}
}
\end{table}

Les variantes à entraînement élargi ou adapté laissent au plus 7 positifs sur un template connu : $2\,874$ des $2\,885$ positifs du test appartiennent à une seule famille de messages (\texttt{ciod: Error reading message prefix on CioStream socket}), apparue peu avant le test. Le protocole G place un entraînement contigu de $80\,000$ lignes juste avant la validation ; il a été conçu après avoir observé la position de cette famille, et utilise donc une information sur le test. $2\,880$ positifs du test portent alors un template connu. Le seuil vaut $6{,}9770$ et la matrice de confusion compte $2\,879$ vrais positifs, 349 faux positifs, $16\,674$ vrais négatifs et 6 faux négatifs. Sur les $19\,723$ événements à template connu, dont $2\,880$ positifs, Histogram atteint une précision de $0{,}9347$, un rappel de $0{,}9983$ et un F1 de $0{,}9654$. La part d'événements de test à template inconnu tombe à $0{,}93\,\%$, et la règle « sévérité absente de l'entraînement » n'obtient plus qu'un F1 inférieur à $0{,}001$ : la décision ne repose plus sur la nouveauté de sévérité.

La couverture du parseur dépend donc surtout de la proximité temporelle entre entraînement et test, les familles d'anomalies de BGL survenant en rafales ; Drain3 reconnaît tous les messages de la famille dominante dès qu'il l'a rencontrée. Le F1 varie de $0{,}5985$ à $0{,}9419$ selon la variante, sans évolution monotone avec la taille de l'entraînement. Le score de G décrit ce que détecte Histogram lorsque cette famille est connue ; il ne constitue pas une performance attendue, et aucune variante ne remplace le protocole principal.
